\setcounter{section}{0}

{\Large\bfseries\filcenter
Mentoring Plan
\\\rule{\textwidth}{0.4pt}}

% No more than 3 pages!!! 
The mentoring plan aims to provide structured support and guidance to one undergraduate and one graduate student involved in this CAREER project. The plan focuses on fostering academic growth, professional development, and personal skills, with a particular emphasis on robotics security, privacy, and emerging threats in that field. The mentorship will involve regular meetings, hands-on project participation, research opportunities, and career preparation and readiness.

\section*{Undergraduate Student Mentoring}

The undergraduate student mentoring aims to enhance their understanding of robotics, security, and related fields, develop basic research skills, prepare for career opportunities in robotics/AI and security, and improve communication, teamwork, and problem-solving abilities.
The student will be introduced to the project, team members, and expectations during the first meeting of joining the project. 
This meeting will also involve discussing the student's interests, goals, and prior experience. Regular one-on-one meetings will be scheduled weekly to discuss progress, challenges, and next steps, with feedback provided on tasks and assignments.
% Hands-on project involvement will be a key part of the mentoring plan. 
The student will be assigned specific tasks such as programming, data analysis, or mechanical assembly, ensuring they have practical experience with robotic systems and tools. 
To complement this, the student will be exposed to research through literature reviews and participation in lab meetings and discussions.
Attendance at relevant workshops, seminars, and lectures will be encouraged to broaden the student's knowledge. 
%These events will cover various topics, including emerging threats and countermeasures in robotics, security, and privacy. Skill development sessions will be organized to enhance both technical and soft skills, with resources provided for self-paced learning.
Career guidance will focus on discussing career paths and opportunities.
Assistance with resume building, interview preparation, and networking will also be provided. 
%Mid- and end-of-term evaluations will be conducted to assess the student's progress and provide constructive feedback for continuous improvement.

\section*{Graduate Student Mentoring Plan}
For the graduate student, the mentoring plan aims to deepen their knowledge in specialized areas of robotics/AI and cybersecurity, develop advanced research skills, prepare for leadership roles in academia or industry, and improve communication, leadership, and project management skills. 
A strong focus will be on understanding and addressing robotics security, privacy, and emerging threats.
An initial meeting will be held to discuss the student's research interests, career goals, and previous experiences, setting clear expectations, and outlining the mentoring plan. 
Bi-weekly meetings will be scheduled to discuss research progress and academic challenges, with detailed feedback provided on research work and publications.
The student will be given leadership roles in specific research projects or components within the project, encouraging the development of research proposals and applications for funding or grants. They will be involved in cutting-edge research, including robotics security and privacy issues, and co-authoring research papers for conferences and journals.
Participation in advanced workshops and seminars will be supported, particularly those that focus on robotics security, privacy, and emerging threats. The student will also be encouraged to present their research findings at conferences and seminars. 
%Skill development sessions will provide opportunities to enhance both technical and soft skills, with resources for continuous learning and professional development.
Career preparation will involve discussions on career paths and opportunities in academia, research institutions, and the industry. Assistance will also be provided with building a professional network and preparing for job applications and interview processes. 
%Regular evaluations will be conducted to assess research progress and skill development, with comprehensive feedback and guidance on academic and professional growth.

